%===============================================================================
% 致谢
%===============================================================================

\chapter*{致\hskip1em 谢}
\addcontentsline{toc}{chapter}{致谢}
\fontsize{12pt}{18pt}\selectfont

历时数载的原子核结构探索之旅，至此行将告一段落。回首在近代物理研究所度过的科研时光，从最初面对复杂能谱数据的迷茫困惑，到逐步推导微观模型、亲历在束物理实验、构建能级纲图并完成本篇学位论文，每一个坚实脚印的背后都凝聚着师长、同侪与家人的无私关怀与支持。

首先，谨向我的导师致以最崇高的敬意与最衷心的感谢。导师严谨求实的科学态度、敏锐深邃的物理洞察力以及对原子核物理前沿的宏阔视野，深深感染并塑造了我的学术生命。在论文选题、实验设计、物理图像提炼以及文稿撰写的全过程中，导师倾注了大量心血，给予了悉心指导与全力支持。

感谢近代物理研究所各加速器运维机组（SFC/SSC/CSR）与在束 $\gamma$ 谱仪实验团队的全体工作人员，感谢你们在繁重的束流运行实验中保障了高质量重离子束流的稳定供给与探测系统的精确运转。

感谢课题组各位师兄、师姐与同学在科研合作、数据离线分析算法探讨以及日常生活中的热情相伴与无私帮助。

感谢我的家人多年来默默的理解、包容与坚强后盾。你们无条件的爱与鼓励，是我在探索未知物理世界道路上永不枯竭的精神动力之源。

最后，向在百忙之中抽时间评审本篇学位论文并参加论文答辩的各位专家教授致以诚挚的谢意！
